\documentclass[10pt,a4paper]{amsart}
\usepackage[margin=19mm]{geometry}
\usepackage{amsmath,amssymb,mathtools}
\usepackage[scr=rsfs,cal=euler]{mathalfa}
\usepackage{xfrac}
\usepackage[unicode,hidelinks,bookmarks=false]{hyperref}
\newcommand{\ie}{i.e.\ }
\newcommand{\cf}{cf.\ }
\newcommand{\resp}{respectively\ }
\newcommand{\Subsection}{\subsection}
\providecommand{\nobreakdash}{\nobreak\hbox{-}}
\setlength{\emergencystretch}{2em}
\multlinegap0pt
\usepackage{bm}
\usepackage{bbm}
\usepackage{dsfont}
\usepackage{xspace}
\usepackage{cancel}
\usepackage{mathtools}
\usepackage{amsthm}
\makeatletter
\@ifundefined{thm}{\newtheorem{thm}{Thm}}{}
\@ifundefined{lem}{\newtheorem{lem}[thm]{Lemma}}{}
\makeatother
\def\cH{{\mathcal{H}}} \def\cI{{\mathcal{I}}} \def\cJ{{\mathcal{J}}} \def\cK{{\mathcal{K}}} \def\cL{{\mathcal{L}}} \def\cM{{\mathcal{M}}} \def\cN{{\mathcal{N}}} \def\cO{{\mathcal{O}}} \def\cP{{\mathcal{P}}} \def\cQ{{\mathcal{Q}}} \def\cR{{\mathcal{R}}} \def\cS{{\mathcal{S}}} \def\cT{{\mathcal{T}}} \def\cU{{\mathcal{U}}} \def\cV{{\mathcal{V}}} \def\cW{{\mathcal{W}}} \def\cX{{\mathcal{X}}} \def\cY{{\mathcal{Y}}} \def\cZ{{\mathcal{Z}}}
\newcommand\Meas{{\operatorname{Meas}}}
\newcommand\Ret{\mathtt{Ret}}
\newcommand{\G}{\Gamma}
\title[Metric criteria for fixed price of countable groups]{Metric criteria for fixed price of countable groups}
\author{Erin Bevilacqua, Lewis Bowen}
\date{Source: arXiv:2510.05459v1 (CC BY 4.0)}
\begin{document}
\maketitle

\noindent\emph{Source and licence.} Metric criteria for fixed price of countable groups. Version arXiv:2510.05459v1 (CC BY 4.0), distributed under the Creative Commons Attribution 4.0 International licence, as recorded in the arXiv licence metadata for this exact version and shown on the abstract page https://arxiv.org/abs/2510.05459v1. Full rights notice: licenses/arxiv-2510.05459-CC-BY-4.0.txt.\par\smallskip

\noindent\textit{Editorial scope.} Counted unit: the Lem whose source label is \texttt{L:half-way} in the article (source0.tex lines 2569--2573, source label \texttt{L:half-way}), delivered together with the article's own complete original proof of it (source0.tex lines 2575--2643, one \texttt{proof} environment of 69 lines). No mathematics is written, repaired or re-derived: statement and proof are the article's own text line for line. Required context retained uncounted: none. Every definition, display and label of those lines is retained unchanged. Omitted from this package and not redistributed: 1--2568, 2574--2574, 2644--3397. Nothing inside a retained excerpt is omitted or altered: every retained body is delivered exactly as the article's lines are written, so no omission receipt of any kind accompanies this package. Citations: the retained excerpts carry no citation command at all, so no bibliography entry of the article is transcribed and no reference is redistributed. Reference adaptation: none was needed and none was made; every reference inside the delivered unit resolves inside it, so no unresolved reference or citation is produced. Printed slips kept literal: no printed word, symbol, hypothesis or number of the counted statement or of its proof was altered. Layout and class: the article's own class is \texttt{article} with packages including graphicx, amsmath, amsthm, amsfonts, amscd, amssymb, comment, eucal, latexsym, mathrsfs, stmaryrd, xy, epsfig, fancyhdr; the wrapper is the lane's pinned \texttt{amsart} editorial wrapper at 10pt on A4, which restates the theorem-like environments the retained text needs and adds the article's own macro definitions that the retained text needs (\texttt{\textbackslash G}, \texttt{\textbackslash Meas}, \texttt{\textbackslash Ret}, \texttt{\textbackslash cH}), with extra wrapper packages (bm, bbm, dsfont, xspace, cancel, mathtools). The delivered numbering is the unit's own. Assets: no retained span contains a figure environment, a graphics-inclusion command, a PDF-page inclusion, a table or a TikZ picture; no image file is copied into this package, the package declares no asset dependency and no image file is redistributed with it. Licence: version arXiv:2510.05459v1 of this article is distributed under the Creative Commons Attribution 4.0 International licence, as recorded in the arXiv licence metadata for this exact version and shown on the abstract page https://arxiv.org/abs/2510.05459v1. A full rights notice is delivered at licenses/arxiv-2510.05459-CC-BY-4.0.txt. Characters: every retained excerpt is pure ASCII, so the delivered engine, XeLaTeX with its default Latin Modern text font, reports no missing character.
\begin{lem}\label{L:half-way}
% Suppose $\mu \in \Meas(\cH)$ satisfies that $\Expand_*\mu$ is equivalent to $\mu$. Let $g\in\G$.
Let $\nu \in \partial\Meas(\cH)$ and $g\in \G$. 
Then for $\nu^2$-a.e.\ $h=(h_1,h_2)\in Z'_g$, $\Ret(Z_g, h)$ is infinite.
\end{lem}

\begin{proof}
It suffices to prove that for every $m>0$ %there exists $r(m)>0$ such that
\begin{align}\label{E:Z'-g}
 \lim_{r\to\infty}   \nu^2(Z'_g \setminus Z'_{g,m,r})=0.
\end{align}
To see this, assume \eqref{E:Z'-g}. Then for all $m>0$, for $\nu^2$-a.e.\ $h \in Z'_g$ there is an $r$ (depending on $m$ and $h$) such that $h\in Z'_{g,m,r} \subset Z'_{g,m}$. Thus $Z'_g \subset \bigcap_{m=1}^\infty  Z'_{g,m}$ (up to measure zero) which implies $\nu^2$-a.e.\ $h\in Z'_g$ has 
infinite return times to $Z_g$: $|\Ret(Z_g, h)|=\infty$.



For the remainder of the proof, let $\nu \in \partial\Meas(\cH)$. Thus there exists a sequence $(\mu_{n_i})_{i=1}^\infty$ with $n_i\to\infty$ such that $\mu_{n_i}$ converges almost-weakly to $\nu$ and each $\mu_{n_i}$ is of the form $\mu_{n_i}=|B(n_i)|^{-1}\sum_{x\in\G} \delta_{d_x-n_i}$.

Both sets $Z'_g$ and $Z'_{g,m,r}$ are closed and open in $\cH$ and both are contained $Z_g = \cH_{<0} \times g\cH_{<C_q}$. Because the restriction of $\mu_{n_i}$ to $\cH_{\le t}$ converges weakly to the restriction of $\nu$ to $\cH_{\le t}$ for any $t$, we must also have that the restriction of $\mu^2_{n_i}$ to $Z'_g$ converges weakly to the restriction of $\nu^2$ to $Z'_g$.   So the Portmanteau Theorem implies for any $m,r>0$
\begin{align}
 \lim_{i\to\infty}  \mu^2_{n_i}(Z'_g \setminus Z'_{g,m,r})= \nu^2(Z'_g \setminus Z'_{g,m,r}).
\end{align}
This equality is one of the main reasons for using the almost-weak topology instead of the vague topology. 

It now suffices to prove 
$$\lim_{r\to\infty} \liminf_{i\to\infty}  \mu^2_{n_i}(Z'_g \setminus Z'_{g,m,r}) = 0.$$

By definition of $\mu_{n}$, 
\begin{align*}
\mu^2_n(  Z'_g \setminus Z'_{g,m,r}) &= \frac{\#\{(x,y) \in B(e,n)\times B(g,n):~(d_x-n,d_y-n) \in Z'_g \setminus Z'_{g,m,r}\}}{|B(n)|^2}.
\end{align*}
This is because if $(d_x-n,d_y-n) \in Z'_g$ then $(x,y) \in B(e,n)\times B(g,n)$:
For the first coordinate it is clear that $x$ must be in $B(n)$ to ensure that $d_x(e)-n = |x| - n<-C-|g|$. 
For the second coordinate, if $d_y(g)-n<-C-2|g|-C_q<0$ then $d(y,g)<n$.

By definition, $(d_x-n,d_y-n) \in Z'_g$ if and only if 
\begin{align*}
    |x| &\le n - C-|g|,\quad 
    d(y,g)  \le n - C - C_q- 2|g|.
    \end{align*}
On the other hand, if $(d_x-n,d_y-n)\in Z'_g$ and $f \in \Ret(Z_g, (d_x-n,d_y-n))$ then $f(d_x-n)(e)\le 0$ and $f(d_y-n)(g)\le C_q$. Equivalently, 
\begin{align*}
d(x,f^{-1})\le n, \quad d(y,f^{-1}g)\le n+C_q.
\end{align*}
This is equivalent to
$$f^{-1} \in B(x,n) \cap B(y,n+C_q)g^{-1}.$$
Letting $N=n-C-|g|$ and $M=n-C-2|g|-C_q$ it suffices to show
\begin{align*}
\lim_{r\to\infty} \liminf_{n\to\infty}  \frac{\#\{(x,y) \in B(e,N)\times B(g,M):~|B(r) \cap B(x,n) \cap B(y,n+C_q)g^{-1}| < m\}}{|B(n)|^2} =0.
\end{align*}
The overlapping neighborhoods property is equivalent to: for every $m>0$
\begin{align*}
\lim_{r\to\infty} \liminf_{n\to\infty}  \frac{\#\{(x,y) \in B(n)^2:~|B(r) \cap B(x,n+C) \cap B(y,n+C)| < m\}}{|B(n)|^2} =0.
\end{align*}
By replacing $n$ with $n-C-|g|$ we see that this is equivalent to:
\begin{align*}
\lim_{r\to\infty} \liminf_{n\to\infty}  \frac{\#\{(x,y) \in B(n-C-|g|)^2:~|B(r) \cap B(x,n-|g|) \cap B(y,n-|g|)| < m\}}{|B(n-C-|g|)|^2} =0.
\end{align*}
We claim that
$$\{(x,y) \in B(e,N)\times B(g,M):~|B(r) \cap B(x,n) \cap B(y,n+C_q)g^{-1}| < m\}$$
is contained in
$$\{(x,y) \in B(n-C-|g|)^2:~|B(r) \cap B(x,n-|g|) \cap B(y,n-|g|)| < m\}.$$
To see this, suppose $(x,y)$ is in the first set. It is immediate that $x\in B(n-C-|g|)$. By the quasi-triangle inequality, $|y| \leq d(y,g)+d(g,e)+C_q\leq n-C-|g|$, so $y\in B(n-C-|g|)$ too. To finish, it suffices to show that
$$B(r) \cap B(x,n-|g|) \cap B(y,n-|g|) \subset B(r) \cap B(x,n) \cap B(y,n+C_q)g^{-1}.$$
This holds because $B(x,n-|g|) \subset B(x,n)$ and $B(y,n-|g|) \subset  B(y,n+C_q)g^{-1}$. To justify the latter inclusion, let $z \in B(y,n-|g|)$. This means $d(y,z)\le n-|g|$. By the quasi-triangle inequality and left-invariance,
$$d(y,zg) \le d(y,z) + d(z,zg)+C_q = d(y,z) + |g|+C_q \le n+C_q.$$
Thus $zg \in B(y,n+C_q)$ which implies $z \in B(y,n+C_q)g^{-1}$ as required.

It follows that 
\begin{align*}
&\lim_{r\to\infty} \liminf_{n\to\infty}  \frac{\#\{(x,y) \in B(e,N)\times B(g,M):~|B(r) \cap B(x,n) \cap B(y,n+C_q)g^{-1}| < m\}}{|B(n)|^2} \\
&\leq \lim_{r\to\infty} \liminf_{n\to\infty} \frac{\#\{(x,y) \in B(n-C-|g|)^2:~|B(r) \cap B(x,n-|g|) \cap B(y,n-|g|)| < m\}}{|B(n-C-|g|)|^2}.
\end{align*}
The latter equals zero by the overlapping neighborhoods property. This completes the proof.
\end{proof}

\end{document}
